\documentclass[11pt]{article}
\usepackage{amsmath,amssymb,amsthm}
\newtheorem{theorem}{Theorem}
\title{Exact Joint-Zero Fibers under All Pole Permutations}
\author{}
\date{}
\begin{document}
\maketitle

\paragraph{Scope.}
This complete written classification audits a proposed source of
many tangent subspaces. It concerns constancy over the full symmetric
group. A structured set of permutations, such as butterfly networks,
need not satisfy its conclusion. It gives no unrestricted
HeadComplexity lower bound and no superconstant compiler. The exact
checker supplements the algebra; the result is not Lean-checked.

Fix $H\geq3$, distinct pole labels $p_1,\ldots,p_H$, residue labels
$a_1,\ldots,a_H$, and one query with denominator offsets $D_h$ and
numerator offsets $b_h$. Assume $D_h\ne p_k$ for all $h,k$. Define
\[
 G(\pi)=\sum_{h=1}^H\frac{a_{\pi(h)}+b_h}{D_h-p_{\pi(h)}}.
 \tag{1}
\]
An additional query bias can turn a constant value into zero, so
constancy is the relevant joint-zero condition.

\begin{theorem}[Full permutation classification]
If $G$ is constant over all permutations, then one of the following
holds:
\begin{enumerate}
\item All $D_h$ are equal and all $b_h$ are equal.
\item At least two $D_h$ differ, there is a polynomial
$Q(X)=AX^2+BX+C$ with $a_k=Q(p_k)$ for every $k$, and
$b_h=-Q(D_h)$ for every $h$.
\end{enumerate}
Both conditions are sufficient for constancy. In particular, if
$H\geq4$ and the residue labels do not agree with a quadratic on the
pole labels, every head must use the same denominator and numerator
offset at this query.
\end{theorem}
\begin{proof}
Write $C_{hk}=(a_k+b_h)/(D_h-p_k)$. Comparing two permutations that
differ by swapping the labels assigned to two rows gives
$C_{hk}+C_{g\ell}=C_{h\ell}+C_{gk}$. Thus
$C_{hk}=u_h+v_k$ for suitable row and column vectors. Conversely, an
additive matrix has a constant permutation sum.

If every $D_h=D$, the difference of rows is
$(b_h-b_g)/(D-p_k)$. It can be constant in $k$ only if $b_h=b_g$,
because at least two pole labels differ.

Otherwise choose $D_h\ne D_g$. Since the difference of these two
rows is constant in $k$, multiplication by
$(D_h-p_k)(D_g-p_k)$ expresses $(D_g-D_h)a_k$ as a polynomial of
degree at most two in $p_k$. This constructs $Q$.
Polynomial division now gives
\[
 C_{hk}=-Ap_k-(AD_h+B)+\frac{R_h}{D_h-p_k},
 \qquad R_h=b_h+Q(D_h).
 \tag{2}
\]
For any pair of rows with distinct offsets, the difference of the two
last fractions in (2) is constant at all $H\geq3$ pole labels.
After clearing denominators, this is a polynomial identity of degree
at most two with at least three distinct roots. Its quadratic
coefficient first makes that constant zero; its remaining coefficients
then give $R_h=R_g=0$. Every other row can be paired with one of
these two rows with a different offset, proving $R_h=0$ for all $h$.
This proves necessity. The displayed division formula proves
sufficiency in the second case; the first case is immediate.
\end{proof}

\paragraph{The quadratic escape is a continuous fiber.}
For arbitrary, continuously varying poles, put
$a_h=Q(p_h)+u_h$ and $b_{hj}=-Q(D_{hj})$. Then the exact identity
\[
 \sum_h\frac{a_h+b_{hj}}{D_{hj}-p_h}
 =\sum_h\frac{u_h}{D_{hj}-p_h}
   -A\sum_h p_h-A\sum_hD_{hj}-HB
 \tag{3}
\]
holds. The last three terms are row and column corrections. At $u=0$
all pole choices lie in the zero fiber after these corrections.
The tangent readout is therefore the constant-numerator separable
reciprocal family, rather than an automatically independent collection
of $H!$ subspaces. Counting the permutation points in this fiber alone
does not prove that its tangent spaces cover all required orthants.

\paragraph{A genuine structured-permutation escape.}
The full symmetric group assumption matters. Let $H=2^d$ and let
$\pi$ be any setting of the usual butterfly network whose layer $k$
switches bit $k$, with layers in increasing order. At any cut $k$,
the map
\[
 h\longmapsto
 (\text{input bits of }h\text{ at positions }\geq k,
  \text{output bits of }\pi(h)\text{ at positions }<k)
 \tag{4}
\]
is a bijection onto the $H$ possible pairs. Before the cut, higher
bits have not changed and the lower wires are permuted within each
higher block. After the cut, lower bits never change. Consequently
every cost that is an arbitrary function of the pair in (4) has a
butterfly-setting-independent sum. Such costs generally are not
row-plus-column additive on the entire permutation matrix.

For example, a leading pole value depending only on the lower output
prefix and a query denominator depending only on the higher input
prefix produce an invariant sum. A later pole digit can expose
switches at the cut. This is an available zero-branch mechanism, but
it does not yet isolate the different lower output prefixes with
admissible exterior kernels. No shattering theorem is inferred from
(4).
\end{document}
